% !TeX program = xelatex
% title: Times New Roman Bold 16 pt; authors/affiliation/keywords: Times New Roman 11 pt;
% body: Aptos 11 pt; Abstract heading: Aptos Display Bold 14 pt.
% bp units match PDF/Word points (1/72 inch).
\RequirePackage{fix-cm}
\documentclass[11pt,a4paper]{article}
\usepackage[margin=2.5cm]{geometry}
\usepackage{fontspec}
\usepackage{microtype}
\usepackage{parskip}
\setmainfont{Aptos}[
  Path=fonts/, Extension=.ttf,
  UprightFont=*, BoldFont=*-Bold,
  ItalicFont=*-Italic, BoldItalicFont=*-Bold-Italic]
\newfontfamily\timesfont{Times}[
  Path=fonts/, Extension=.TTF,
  UprightFont=*, BoldFont=Timesbd,
  ItalicFont=Timesi, BoldItalicFont=Timesbi]
\newfontfamily\abstractheadingfont{Aptos-Display-Bold.ttf}[Path=fonts/]
\renewcommand{\normalsize}{\fontsize{11bp}{14.5bp}\selectfont}
\setlength{\parindent}{0pt}
\setlength{\parskip}{6bp}
\setlength{\emergencystretch}{1em}
\pagestyle{empty}

\makeatletter
\renewcommand{\maketitle}{%
  \begin{center}
    {\timesfont\fontsize{16bp}{20bp}\selectfont\@title\par}
  \end{center}
  \vspace{12bp}
  {\timesfont\fontsize{11bp}{14bp}\selectfont\noindent\@author\par}
  \vspace{14bp}
}
\makeatother
\renewenvironment{abstract}{%
  \par{\abstractheadingfont\fontsize{14bp}{17bp}\selectfont Abstract\par}
  \vspace{2bp}\normalfont\normalsize
}{\par}

\title{\textbf{Comparing Reinforcement Learning and Model-based Optimization for Energy Management of a Renewable Industrial System}}
\author{Ashish Thapa\textsuperscript{\fontsize{7bp}{8bp}\selectfont 1}, Loukas Kyriakidis\textsuperscript{\fontsize{7bp}{8bp}\selectfont 1}, Ralf Wunderlich\textsuperscript{\fontsize{7bp}{8bp}\selectfont 2}, Eric Pilling\textsuperscript{\fontsize{7bp}{8bp}\selectfont 2}\\[0.3em]
\textsuperscript{\fontsize{7bp}{8bp}\selectfont 1}German Aerospace Center, Institute of Low-Carbon Industrial Processes, Cottbus, Germany\\[0.3em]
\textsuperscript{\fontsize{7bp}{8bp}\selectfont 2}Brandenburg University of Technology, Chair of Economathematics, Cottbus, Germany}
\date{}

\begin{document}
\maketitle

\begin{abstract}
Industrial steam generation is an energy-intensive process, and its decarbonization requires coordinated operation of renewable electricity, high-temperature heat pumps, and thermal energy storage under fluctuating electricity prices, renewable availability, and process demand. The resulting setpoint optimization problem is nonlinear, nonconvex, and constrained, making repeated online model-based optimization computationally demanding. Reinforcement learning (RL) enables direct setpoint decisions after offline training, offering a potential alternative to repeated online model-based optimization. This work investigates RL and model-based optimization for energy management of renewable industrial steam generation. A system comprising a high-temperature heat pump, an on-site wind turbine, grid electricity, and thermal energy storage is considered. The heat pump supplies process steam through the storage system to meet the demand. The energy management problem involves determining heat-pump operation, storage charging and discharging, and grid-wind electricity allocation over a finite horizon. The RL formulation represents the system as a continuous-action Markov decision process, with observations including storage state, operating conditions, electricity prices, and forecasts of wind generation and steam demand. Dependent variables are derived from process equations to reduce the policy search space, while state-dependent action mappings and penalty terms enforce operating constraints. As an optimization benchmark, the same problem is formulated as a constrained nonlinear program and solved using the novel hybrid optimizer “BO-IPOPT”, combining Bayesian Optimization for global exploration with the Interior Point OPTimizer (IPOPT) for local refinement. To account for forecast uncertainty, the optimization benchmark is implemented in a rolling-horizon framework, in which the optimization problem is repeatedly solved over a finite prediction horizon and the resulting decisions are updated as new forecasts and measurements become available. The two approaches are evaluated in terms of accuracy, computational time, constraint satisfaction, and robustness. Last, the use of the RL policy as an initialization for the BO-IPOPT rolling-horizon optimization is investigated. The RL-generated setpoints are provided as an initial solution to BO-IPOPT and compared with a cold-start implementation to assess their impact on solver iterations and computational time while maintaining solution quality and constraint satisfaction. Overall, the study evaluates the potential of RL both as a standalone approach for fast energy management and as a means of accelerating model-based optimization through informed initialization, with the aim of enabling computationally efficient and reliable energy management of industrial utility systems.
\end{abstract}

\begingroup\timesfont\fontsize{11bp}{14bp}\selectfont
\noindent\textbf{Keywords:} Energy management; Reinforcement learning; Model-based optimization; BO-IPOPT; Industrial decarbonization

\par\endgroup

\end{document}
